\documentclass[conference]{IEEEtran}

% ============================================================
% Corrected final-campaign manuscript.
%
% Numerical results are preserved from the supplied final
% manuscript and registered outcome tables.
%
% The former final-results.tex fragment is embedded below.
% Do NOT also add \input{final-results.tex}.
%
% Compile this document separately from supplement.tex.
% The original docs/paper-final/figures/ directory is kept
% unchanged. The six results figures in figures/ here are the
% redesigned versions produced from the same audited rows by
% scripts/plot_final_research_improved.py.
% ============================================================

\usepackage[T1]{fontenc}
\usepackage[utf8]{inputenc}

\usepackage{cite}
\usepackage{amsmath,amssymb}
\usepackage{graphicx}
\usepackage{booktabs}
\usepackage{array}
\usepackage{tabularx}
\usepackage{algorithmic}
\usepackage{placeins}
\usepackage{url}
\usepackage[hidelinks]{hyperref}

\hypersetup{
    pdftitle={
        Predictive Multi-Contact Security Lifecycle Management
        for LEO Spacecraft-Ground Communications
    },
    pdfauthor={Mukhammed Zhumali},
    pdfsubject={
        Final registered ContactAwareSS campaign:
        lifecycle planning, computation, and security consequences
    }
}

\newcommand{\op}[1]{\texttt{#1}}
\newcommand{\missing}{\textit{I}}
\newcommand{\authorcheck}[1]{\textbf{[VERIFY: #1]}}

\newcolumntype{Y}{>{\raggedright\arraybackslash}X}

\setlength{\emergencystretch}{2em}

% Float controls. These do not change IEEE page margins.
\setcounter{topnumber}{3}
\setcounter{dbltopnumber}{3}
\setcounter{totalnumber}{5}

\renewcommand{\topfraction}{0.90}
\renewcommand{\dbltopfraction}{0.90}
\renewcommand{\textfraction}{0.08}
\renewcommand{\floatpagefraction}{0.75}
\renewcommand{\dblfloatpagefraction}{0.75}

% Usage:
% \widefigure{filename.pdf}{maximum height}{caption}{label}
%
% Figures are bounded in height and flushed before Discussion,
% preventing the results figures from appearing after References.
\newcommand{\widefigure}[4]{%
    \begin{figure*}[!tp]
        \centering
        \includegraphics[
            width=0.97\textwidth,
            height=#2,
            keepaspectratio
        ]{figures/#1}
        \caption{#3}
        \label{#4}
    \end{figure*}
}

% ============================================================
% AUTHOR CHECKS
%
% Replace these three values only after checking the frozen
% final-campaign model JSON and environment manifest.
%
% Do not infer them from older SE2/R1/PC2 configurations.
% ============================================================

\newcommand{\CoreMaterialAge}{%
    \authorcheck{final-model value}%
}

\newcommand{\CoreReuseGap}{%
    \authorcheck{final-model value}%
}

\newcommand{\FinalHostHardware}{%
    \authorcheck{final host CPU, RAM, and OS}%
}

\begin{document}

\title{
    Predictive Multi-Contact Security Lifecycle Management
    for LEO Spacecraft--Ground Communications
}

\author{
    \IEEEauthorblockN{Mukhammed Zhumali}
    \IEEEauthorblockA{262937@astanait.edu.kz}
}

\maketitle

% ============================================================
% ABSTRACT
% ============================================================

\begin{abstract}
Time-limited spacecraft--ground contacts constrain both payload
delivery and traffic-key and security-association management.
We study when additional contact visibility changes lifecycle
decisions and whether its benefit justifies computation.
Matched Next and Multi policies share actions, constraints,
costs, a five-component objective, and a full-plan tie rule;
only their visible horizon differs.
A finite-model exact reference is separated from practical
execution using feasible seeds, retained continuations, and
bounded search.
One registered campaign covers geometric ISS, NOAA-19, and HST
schedules, three demand levels, selected sensitivities,
traffic-key disclosure, and channel interruptions.
Of 108 attempts, 93 finish and 15 reach a process watchdog.
All four completed primary horizon pairs tie on the five
objective components; five primary pairs remain incomplete.
No-search controls finish every primary workload.
A prepared exact short-return example retains its additional
served contact, but the bounded autonomous controller does
not recover that benefit.
One late-detection case yields an extra confidential frame
for Multi, while measured computation reduces its useful
transmission time in another case.
The tested geometric core therefore establishes no practical
quality advantage from the additional horizon at the fixed
budget.
Synthetic costs, a logical ground coordinator, and protected
management credentials limit the interpretation.
No onboard validation or general cryptographic-security
guarantee is claimed.
\end{abstract}

\begin{IEEEkeywords}
LEO communications, key lifecycle, security association,
contact-aware planning, bounded search, disruption,
receding horizon
\end{IEEEkeywords}

% ============================================================
% I. INTRODUCTION
% ============================================================

\section{Introduction}
\label{sec:introduction}

A return contact can be long enough to transmit its required
payload but too short to install, activate, and use a replacement
key.
If that requirement is visible during an earlier contact with
the same ground station, some management work can be moved
forward.
The scheduling opportunity depends on key validity, resource
contention, payload requirements, and the selected plan.
A longer forecast alone guarantees neither better service nor
lower cost.

The research question is therefore conditional:
\emph{when does additional contact visibility change joint
traffic-key and Security Association (SA) decisions, and what
service benefit, management cost, and computational burden
follow?}
Multi-contact denotes observation depth for one spacecraft and
several ground stations.
It does not denote a multi-spacecraft optimization problem.

The contribution is a finite lifecycle-planning formulation and
an executable comparison that separates forecast depth from
other policy choices.
The primary objective is admissible service, not compromise
probability.
Disclosure experiments examine consequences of the selected
schedules; they neither introduce advance knowledge of an
attack nor optimize recovery latency.

The framework separates an exact finite-model reference from
a practical policy that can use a feasible but uncertified
incumbent or a common fallback.
It also represents decision delay and actual channel
unavailability explicitly.
These mechanisms expose applicability assumptions without
turning synthetic timing into measurements of a flight
implementation.

The final comparison is interpreted within its registered
workload, demand, incident, and budget settings.
Earlier revisions retain separate evidence identities.
Receding-horizon optimization, contact prediction, advance
authentication, and established cryptographic mechanisms are
not claimed inventions.

% ============================================================
% II. RELATED WORK
% ============================================================

\section{Related Work and Novelty Boundary}
\label{sec:related}

CCSDS SDLS specifies data-link protection, and its Extended
Procedures define key and SA management services
\cite{sdls,ep}.
The Symmetric Key Management Recommended Practice supplies
an explicit key lifecycle and abstract management procedures
\cite{skm}.
Lifecycle states, key distribution, activation, and retirement
therefore pre-exist this study.
We schedule abstract counterparts under a chosen resource
and service model; sharing state names does not establish
SDLS conformance.

Connectivity already informs security-state management.
The IOTA Identity architecture updates an onboard authorization
allowlist during ground contacts and considers revocation delay
and short-lived capability tokens \cite{iota}.
Its object is IoT-device identity and access authorization.
Our object is directional traffic-key and SA timing.
Contact-dependent authorization updates are therefore not
a novelty claimed here.

Delay-Tolerant Networking provides another relevant line of
work.
BPv7 supports bundle delivery in challenged networks
\cite{bp}, while BPSec provides bundle integrity and
confidentiality and leaves key management to separate network
management \cite{bpsec}.
Schedule-Aware Bundle Routing explicitly uses scheduled
contacts for route computation \cite{sabr}.
Prediction and contact-based optimization alone are consequently
insufficient novelty claims.

Most directly, Detke et al.\ couple time-feasible Contact Graph
Routing routes with BPSec key provisioning \cite{detke}.
Their offline mixed-integer program selects routes, active key
scopes, and covered pairs to minimize scopes at a target
connectivity ratio.
Its displayed variables have route, scope, and pair indices;
they do not schedule lifecycle-operation start times.
The paper identifies key scheduling over contact plans as
future work.
Our comparison concerns sequential per-station traffic-key
and SA actions and a controlled visible-horizon change for
a fixed workload.

The control pattern is also established.
Receding-horizon control repeatedly solves a prediction-based
constrained optimization problem and applies an initial
control action before replanning \cite{rhc}.
Here, this pattern is applied to a finite event-based
key/SA lifecycle model.
Results for other optimization classes do not establish the
optimality or real-time feasibility of our implementation.

Table~\ref{tab:related} summarizes the scope of the comparison.
This is a bounded comparison with the listed primary sources,
not an exhaustive review or a claim of priority over all
temporal key-management formulations.
No external authentication or routing protocol has been
reimplemented or empirically outperformed in this study.

\begin{table*}[!t]
    \centering
    \footnotesize
    \setlength{\tabcolsep}{4pt}
    \renewcommand{\arraystretch}{1.15}

    \caption{
        Closest established mechanisms and the bounded comparison
        made here. These are conceptual distinctions, not results
        against reimplementations of the cited protocols.
    }
    \label{tab:related}

    \begin{tabularx}{\textwidth}{
        @{}
        p{0.18\textwidth}
        p{0.30\textwidth}
        Y
        @{}
    }
        \toprule
        Source
        & Established object or mechanism
        & Relationship to this study \\
        \midrule

        SDLS/EP and symmetric key management
        \cite{sdls,ep,skm}
        &
        Space-link protection; key lifecycle; concrete and
        abstract management procedures.
        &
        Provide mechanism and state provenance.
        This study schedules abstract operations under explicit
        costs and constraints, without claiming standards
        conformance.
        \\

        IOTA Identity \cite{iota}
        &
        Device credentials, contact-dependent allowlists,
        and revocation handling.
        &
        Establishes connectivity-aware authorization.
        This study concerns channel traffic keys and paired SAs.
        \\

        BPv7, BPSec, and SABR \cite{bp,bpsec,sabr}
        &
        Bundle delivery, bundle security, and routing over
        scheduled contacts.
        &
        Establish intermittent-connectivity processing and
        contact-aware routing.
        This study holds the selected link workload fixed.
        \\

        Detke et al.\ \cite{detke}
        &
        Time-feasible route catalog coupled to offline
        BPSec key-scope provisioning.
        &
        Closest reviewed provisioning-optimization precedent.
        This study varies observation depth for sequential
        lifecycle execution.
        \\

        Receding-horizon control \cite{rhc}
        &
        Repeated constrained planning followed by execution
        of an initial control action.
        &
        Supplies the established control pattern.
        The contribution is the lifecycle formulation and its
        conditional experimental evaluation.
        \\

        \bottomrule
    \end{tabularx}
\end{table*}

\FloatBarrier

% ============================================================
% III. SYSTEM MODEL
% ============================================================

\section{System and Lifecycle Model}
\label{sec:model}

\subsection{Security State and Validity}

One spacecraft maintains an authorized context for each
configured ground station (GS).
Uplink and downlink use separate traffic keys and SAs;
the two SAs are managed as a pair.
At most one ground-link context is operational.
Station visibility does not imply authorization, and keys
are not silently reused across station contexts.

A contact is represented by its start, end, and station:
\[
    C_i=(s_i,e_i,g_i),
    \qquad s_i<e_i.
\]
Its availability interval is half-open, $[s_i,e_i)$.

The implementation stores immutable metadata and state,
rather than secret key bytes or cipher execution.
Preparation, confirmed installation, activation/binding,
and opening a session are distinct operations.
A prepared key does not authenticate a peer.
Closing a session preserves valid keyed state and persistent
counter progress.
Retired identities are not recycled.
Logical retirement is not evidence of physical erasure.

For creation time $t_k^c$, first activation $t_k^a$,
material-age bound $A_{\max}$, activation cryptoperiod
$C_{\max}$, and committed not-after bound $\bar d_k$,
the active-key deadline is
\begin{equation}
    d_k=
    \min\left\{
        \bar d_k,\,
        t_k^c+A_{\max},\,
        t_k^a+C_{\max}
    \right\}.
    \label{eq:deadline}
\end{equation}

The activation term does not apply before activation.
Reopening or reusing an SA does not extend an existing
deadline.
Between-contact reuse also requires that the gap from the
last successful directional session does not exceed
$G_{\max}$, together with authorization, confirmed
installation, peer synchronization, and counter headroom.

Online confirmations must finish strictly before contact
closure and required-key expiry.
A complete payload frame may end exactly at either boundary.
Expiry is an actual execution event.
If one directional key expires, the common response stops
the paired context and retires its companion according to
the model.
Independently staged material is retained when admissible.

\subsection{Execution and Costs}

The executor supports
\op{PREPARE},
\op{INSTALL},
\op{SWITCH},
\op{OPEN},
\op{REUSE},
one-frame \op{CONTINUE}, and
\op{WAIT}.
Stop and expiry are local transitions.
Incident response is common to all policies.

Operations reserve ground CPU, spacecraft CPU, radio direction,
and waiting phases.
Only executed work is charged, including partial and ultimately
wasted work.
Useful time and delivered service count complete frames.

\begin{table}[!t]
    \centering
    \footnotesize
    \setlength{\tabcolsep}{3.2pt}
    \renewcommand{\arraystretch}{1.12}

    \caption{
        Base synthetic operation profile.
        GS and SC denote ground station and spacecraft.
        These are model assumptions, not hardware measurements.
    }
    \label{tab:costs}

    \begin{tabular}{@{}lrrrrr@{}}
        \toprule
        Operation
        & Elapsed
        & GS CPU
        & SC CPU
        & UL
        & DL \\
        & s & s & s & B & B \\
        \midrule

        PREPARE
        & 2.0 & 0.20 & 0.00 & 0 & 0 \\

        INSTALL
        & 4.0 & 0.10 & 0.20 & 600 & 200 \\

        SWITCH
        & 1.0 & 0.05 & 0.10 & 100 & 100 \\

        OPEN/REUSE
        & 1.0 & 0.02 & 0.04 & 40 & 40 \\

        Revocation notice
        & 0.4 & 0.02 & 0.02 & 64 & 32 \\

        \bottomrule
    \end{tabular}
\end{table}

The profile in Table~\ref{tab:costs} models a 400~B/s
management channel.
Elapsed duration includes waiting and is not simply the
sum of CPU measurements.
PREPARE uses ground resources only and may overlap payload
if reservations permit.
Online management competes for the link.
Stop and expiry have zero modeled cost in this profile.

The state includes pending operations, committed frames,
per-contact delivery, and accrued costs.
Replanning does not cancel committed work or charge it again.
Confirmed state changes occur only at operation checkpoints.
An uncertain SWITCH, OPEN, or REUSE remains blocked pending
explicit reconciliation.
Restoring connectivity alone does not reconcile the endpoints.

\subsection{Deployment Interpretation}

The model assumes a \emph{logical ground coordinator}
with access to the permitted forecast and a consistent abstract
execution state.
Ground-only preparation can occur while disconnected.
Installation, switching, opening, reuse, and spacecraft notice
require the matching available link.
Local stop and expiry do not require successful delivery
of a new command.

This identifies the planning location and timing boundary,
but it is not a complete deployment model.
Trusted management credentials and abstract confirmations
are assumed.
Distributed telemetry/state reconciliation, an onboard command
protocol, controller failover, and measured command/telemetry
bytes are not implemented.

The single logical state is an explicit abstraction,
particularly when a downlink continues after ground
revocation notice.
A flight realization must map the abstract operations to a
concrete protocol and validate their costs and information flow.

\subsection{Conditional Installation Opportunity}
\label{sec:opportunity}

For two specified return-entry states, define
\begin{align}
    \tau_{\mathrm{installed}}
        &=t_S+t_O+t_P,
        \label{eq:installed-chain}\\
    \tau_{\mathrm{prepared}}
        &=t_I+\tau_{\mathrm{installed}}.
        \label{eq:prepared-chain}
\end{align}

Here $t_I$, $t_S$, and $t_O$ are installation, switching,
and opening durations.
The quantity $t_P$ is the time needed for the required
complete-frame payload, and $W$ is the remaining contact time.

If
\begin{equation}
    \tau_{\mathrm{installed}}
    \leq W
    <
    \tau_{\mathrm{prepared}},
    \label{eq:opportunity}
\end{equation}
only the installed-entry chain fits.

This is a feasibility statement about these two chains.
All validity, counter, resource, and boundary checks must
still hold.
It is not a lower bound over every policy or initial state,
and it does not assert that the final practical campaign
realizes a service advantage.

% ============================================================
% IV. POLICY AND PLANNING
% ============================================================

\section{Horizon, Objective, and Policy Modes}
\label{sec:planning}

\subsection{Observation and Objective}

For a decision inside contact $C_i$, $H$ counts
\emph{additional future contacts}:
\begin{equation}
    \mathcal{W}_i^H=
    \left(
        C_i,C_{i+1},\ldots,C_{\min(n,i+H)}
    \right).
    \label{eq:horizon}
\end{equation}

Next uses $H=1$ and sees the current and next contact.
The basic Multi comparison uses $H=2$ and sees the current
contact and two future contacts.

Between contacts there is no current contact, so the nearest
$H$ future windows are visible.
A search freezes this observation and its endpoint.
Advancing a hypothetical branch does not reveal another contact.

Future prescribed actions, channel outages, restoration times,
attack occurrences, and unobserved failures are excluded.
Nominal branches model future confirmations as successful
until actual execution supplies contrary evidence.
Only the first selected action is admitted before the
controller observes and replans.

Under the common admission constraints, plans are compared
using the lexicographic objective
\begin{equation}
    \max_{\mathrm{lex}} J
    =
    \left(
        N_{\mathrm{served}},
        T_{\mathrm{useful}},
        -C_{\mathrm{CPU}},
        -B_{\mathrm{mgmt}},
        -N_{\mathrm{new}}
    \right).
    \label{eq:objective}
\end{equation}

Service requires the configured complete-frame minimum
separately in each direction.
CPU and management bytes sum both endpoints or directions.
New keys count directional identities, including keys produced
by wasted preparation.

Previously accrued work is retained and already served
contacts are not counted twice.
There is no terminal bonus for unused fresh material.
Recovery latency, readiness, and compromise probability are
not additional objectives.

Only equality of all five components invokes a full-plan
tie rule.

CANONICAL orders preparation, installation, switching,
opening, reuse, continuation, and wait, followed by start time
and the complete identity/end-time tuple.
Plans are compared at their first difference, with a proper
prefix preceding its extension.

CONTINUE\_WAIT places continuation first, then wait, then
management actions.
Within a group it compares earlier end time before the original
tuple.
It is therefore more specific than a permutation of action names.

Search traversal order is distinct from the full-plan preference.
Within each rule and execution mode, Next and Multi share
all policy mechanisms except the permitted horizon.

\subsection{Exact Reference}

The exact reference enumerates a finite event-based action
model.
Candidate starts are discrete.
At most one spare pair or preparation is admitted per context.
Frame direction follows the common executor, and actions
must satisfy completion-within-horizon restrictions.

A root status \op{COMPLETE} certifies the finite continuation
within this model.
A root status \op{INCOMPLETE} does not.
The reference controller stops on an incomplete search and
preserves a checkpoint instead of executing an unproved fallback.

Neither completed root searches nor certified suffix reuse
prove global optimality of the executed receding-horizon policy.

\subsection{Practical Candidate Portfolio}
\label{sec:portfolio}

The practical policy is separate from the exact reference.
A physical execution marked \op{FINISHED} can contain
incomplete searches.
It must not be interpreted as a sequence of proven optima.

The final practical controller compares three sources of
fully evaluated feasible continuations:

\begin{enumerate}
    \item a deterministic feasible seed;
    \item a re-evaluated retained continuation;
    \item terminal candidates obtained by bounded search.
\end{enumerate}

The seed is generated by the shared feasible-plan builder,
using only the permitted observation and common admission
and execution mechanisms.
The no-DFS controls use this same construction without
depth-first search.

Retention can cross completed preparation, installation,
switching, opening, and frame steps.
Executed costs are not charged again.
If the visible horizon extends, the retained suffix is
completed using the same seed procedure.
Stale, inconsistent, or inadmissible continuations are rejected.

Every candidate is evaluated from the same actual observable
root to the same current endpoint.
Old scores from different roots or endpoints are not compared
directly.
The seed and retention mechanisms receive neither hidden
contacts nor future incidents.

All available complete continuations compete under the
unchanged five-component objective and full-plan comparator.
There is no special preference for retention.

For the set $\mathcal{P}_{\mathrm{eval}}$ of available,
fully evaluated, admissible candidates at one root,
the completed selection satisfies
\begin{equation}
    J(\pi_{\mathrm{selected}})
    \succeq_{\mathrm{lex}}
    J(\pi),
    \qquad
    \forall \pi\in\mathcal{P}_{\mathrm{eval}}.
    \label{eq:portfolio-guarantee}
\end{equation}

This is a local available-candidate guarantee.
It does not imply mission-wide dominance, global optimality,
or superiority of Multi over Next.
The source search remains incomplete when its search budget
was exhausted.
Strictly certified suffix reuse is a narrower mechanism with
its own verification contract.

Figure~\ref{fig:controller} summarizes the execution boundary.
The pseudocode describes the portfolio and dispatch contract;
it does not replace the frozen implementation's action
generation or candidate-ordering rules.

\begin{figure}[!t]
    \centering
    \begin{minipage}{0.98\columnwidth}
        \hrule
        \vspace{4pt}
        \footnotesize

        \textbf{At each actual decision boundary}

        \begin{algorithmic}[1]
            \STATE Process mandatory physical events and
                   preserve committed work.

            \IF{an observed incident requires emergency handling}
                \STATE Execute the next admissible common
                       response step and return to the event loop.
            \ENDIF

            \STATE Freeze the permitted observation $O_H$,
                   current root, and evaluation endpoint.

            \STATE Start the common decision-time budget.

            \STATE Build and evaluate a feasible seed.

            \STATE Replay and validate the retained suffix;
                   extend it to the current endpoint when needed.

            \STATE Run bounded search within the remaining
                   registered limits.

            \STATE Collect available fully evaluated feasible
                   candidates from all three sources.

            \IF{the decision deadline expired}
                \STATE Select the common admissible safe response.
            \ELSIF{a feasible candidate is available}
                \STATE Select the best candidate using
                       Eq.~\eqref{eq:objective} and the full-plan
                       tie rule.
            \ELSE
                \STATE Use the common admissible fallback.
            \ENDIF

            \STATE In measured mode, advance actual events over
                   the charged decision delay.

            \STATE Revalidate the first intended action against
                   the resulting actual state.

            \IF{the intent is stale or its WAIT event was
                reached or changed}
                \STATE Request a fresh, separately charged decision.
            \ELSE
                \STATE Execute only the first admitted action;
                       retain its feasible continuation.
            \ENDIF

            \STATE Return to the physical event loop.
        \end{algorithmic}

        \vspace{4pt}
        \hrule
    \end{minipage}

    \caption{
        Practical receding-horizon controller.
        Seed construction, suffix replay, search, and selection
        share the decision budget.
        A feasible practical selection does not certify an
        incomplete search.
    }
    \label{fig:controller}
\end{figure}

\subsection{Search Budgets and Operational Baselines}

The final campaign fixes 5000 nodes and 1~s per search,
2~s per complete decision, and a separate 600~s process
watchdog.
Seed and suffix stages each have a 350~ms sublimit and
a 10\,000-transition cap.
Their work, search, and selection all consume the common
decision budget.

There is no additional cumulative controller-time cutoff.
One JVM with a 1~GiB heap runs at a time.
The limits are not raised after observing incomplete runs.
Host budgets are separate from modeled security-operation CPU.

FixedPeriod is an author-defined operational baseline,
not a mandated CCSDS scheduler.
Its renewal timer is anchored to the last successful SWITCH.
Overdue requests coalesce, and a missed renewal does not
extend key validity.
One selected period cannot establish superiority over every
possible periodic policy.

PerContact requires newly activated material for the current
contact under its registered preparation rule.
Both operational controls share admission checks and the
event executor.
The primary attribution of a horizon effect remains the
matched Next/Multi comparison.

\FloatBarrier

% ============================================================
% V. DELAY AND DISRUPTION
% ============================================================

\section{Delay and Channel Availability}
\label{sec:availability}

\subsection{Decision and Dispatch Delay}

Offline computation does not advance physical simulation time.
Measured mode introduces the critical-path delay
\begin{equation}
    \delta=
    \delta_{\mathrm{fixed}}+\alpha t_{\mathrm{host}},
    \qquad
    \delta_{\mathrm{fixed}},\alpha\geq 0.
    \label{eq:delay}
\end{equation}

Here $t_{\mathrm{host}}$ includes observation and candidate
construction, seed, replay, search, and selection.
Checkpoint persistence and physical execution are measured
separately and are not included in this modeled delay.

The final measured block uses
$\delta_{\mathrm{fixed}}=0$ and $\alpha=1$.
This is a host-dependent critical-path assumption, not a
measured spacecraft deployment.
The coordinator issues no new action while computing,
while committed work and physical events continue.

After the delay, the controller constructs a fresh observation
and revalidates the first selected intent.
Reached or changed WAIT targets trigger another measured
decision.
An expired intent is not silently retargeted to contact closure.
Other stale intents similarly require replanning after
physical time advances.

Decision deadline misses use the common admissible wait
fallback.
An observed incident supersedes the scheduled choice through
the common response.

Shifted execution has no inherited optimality certificate.
Aggregate host runtime cannot simply be subtracted from a
service or recovery benefit.
Only delay on the relevant critical path has that meaning,
and no onboard CPU equivalence is inferred from $\alpha$.

\subsection{Controlled Outages Without Future Leakage}

An interruption names a contact and an actual unavailable
interval wholly inside the original forecast window.
Full-window loss represents a missed contact.
Initial loss represents delayed availability, and trailing
loss represents early closure.
The original forecast remains unchanged.

Actual DOWN aborts a committed frame, fails an online
primitive with its consumed cost, and stops operational use
locally.
PREPARE remains possible because it uses ground resources.

Actual UP records availability but does not automatically
reopen an SA, confirm an interrupted operation, or clear
uncertainty.
INSTALL can be retried when subsequently admissible.
Uncertain SWITCH, OPEN, or REUSE remains fail-closed because
reconciliation is not implemented.

The observation first selects the original $H$-limited prefix
and then removes a currently unavailable contact using only
past DOWN/UP events.
It does not reveal a farther contact to replace the unavailable
one.
Future restoration events are not exposed to the planning
branch.
Actual UP wakes execution and permits a new observation.

This is conservative handling of known loss, not estimation
of its duration or a complete stochastic model of forecast error.

Frame completion precedes loss at the same timestamp.
Expiry and original contact closure precede online confirmation.
Loss precedes an acknowledgement at the same instant.
UP precedes DOWN for adjacent outage intervals, avoiding
a fictitious available instant.
These are explicit simulation conventions.

\subsection{Aborted-Frame Accounting}

Each aborted committed frame increments its key counter once
and updates any still-bound SA mirror.
It contributes zero delivered frames, delivered bytes,
and useful time.
Even a zero-byte prefix conservatively consumes its nonce.

This prevents a temporarily interrupted context from reopening
with counter progress that omits an admitted transmission.
The operation accounts for an already admitted frame;
it does not grant new permission to transmit.

A ground-rejected committed downlink can consume its retired
key's counter without resurrecting that key or refreshing
successful-use history.
Counter audits therefore include delivered and aborted
committed frames.

Historical files retain their executable identities.
This behavior is not retroactively attributed to earlier
implementations.

% ============================================================
% VI. EXPERIMENTAL DESIGN
% ============================================================

\section{Experimental Design and Metrics}
\label{sec:experiments}

\subsection{Public Data and Geometric Workloads}

The geometric workloads use archived CelesTrak
general-perturbations orbital elements \cite{celestrak}
and publicly documented ground-station coordinates obtained
from SatNOGS observation records \cite{satnogs}.
Orekit supplies the orbital and visibility calculations
\cite{orekit}.

A common selector, independent of the security policies,
chooses whole nonoverlapping windows under the single-link
assumption.
The workloads are geometric opportunities, not records of
authorized protected spacecraft--ground exchanges.
Radio-link conditions, antenna availability, and operational
authorization are not inferred from visibility.

The base schedules contain 46 ISS, 38 NOAA-19, and 48 HST
contacts.
ISS uses the saved October 10, 2026 day.
NOAA-19 and HST use October 25, 2026.
Saved GP epochs are October 10 for ISS and October 24
for NOAA-19 and HST.

All three use the same 32-station cohort, a $10^\circ$
elevation threshold, and zero switching guard.
The dates and source snapshots are preserved:
geometry is not refreshed after observing outcomes.

\subsection{Controlled X5 and Y Scenarios}

The X5 and Y regressions retain the original three-contact
geometry summarized in Table~\ref{tab:witness}.
Both initially contain valid installed contexts for A and B.
The initial A material expires at $t=70$~s, while initial B
material expires at $t=120$~s.

The autonomous regression selects and pays for any new
preparation itself.
The prepared exact reference is a separate fixture with a
prescribed common PREPARE before the contacts.
Its preparation cost is included exactly once.
These two initialization conventions are not pooled into one
result.

The original witness uses 100~B complete frames at 400~B/s
and requires 400~B in each direction to serve a contact.
This corresponds to 2~s of minimum useful transmission over
the shared alternating link.

\begin{table}[!t]
    \centering
    \footnotesize
    \setlength{\tabcolsep}{5pt}
    \renewcommand{\arraystretch}{1.15}

    \caption{
        Controlled contact geometry.
        Times are seconds relative to the start of $C_1$.
    }
    \label{tab:witness}

    \begin{tabular}{@{}llll@{}}
        \toprule
        Scenario & $C_1$ & $C_2$ & $C_3$ \\
        \midrule
        X5 & A: $[0,20)$ & B: $[30,50)$ & A: $[80,85)$ \\
        Y  & A: $[0,20)$ & B: $[30,50)$ & B: $[80,85)$ \\
        \bottomrule
    \end{tabular}
\end{table}

\subsection{Registered Matrix and Parameters}

The immutable final registration contains 108 attempts:
54 complete-schedule offline configurations,
15 measured ISS configurations,
18 historical-mechanism regressions,
nine sensitivity cases,
ten disclosure/channel cases, and
two management-interruption cases.

Table~\ref{tab:parameters} gathers the main final-campaign
settings.
All cost and workload assumptions are retained from the
registered models, rather than adjusted to improve the outcome
of a selected policy.

\begin{table*}[!t]
    \centering
    \footnotesize
    \setlength{\tabcolsep}{4pt}
    \renewcommand{\arraystretch}{1.12}

    \caption{
        Final-campaign parameters and their interpretation.
        VERIFY fields must be copied from the frozen final-model
        and environment records before submission.
    }
    \label{tab:parameters}

    \begin{tabularx}{\textwidth}{
        @{}
        p{0.24\textwidth}
        p{0.29\textwidth}
        Y
        @{}
    }
        \toprule
        Parameter & Value & Interpretation \\
        \midrule

        Selected geometric schedules
        &
        ISS: 46; NOAA-19: 38; HST: 48 contacts
        &
        Full selected sequences; not protected-link measurements.
        \\

        Station cohort and geometry
        &
        32 stations; $10^\circ$ mask; 0~s guard
        &
        Common geometry settings; station authorization is
        a separate model assumption.
        \\

        Primary cryptoperiod $C_{\max}$
        &
        5400~s
        &
        Maximum use after first activation.
        Sensitivity uses 2700 and 10\,800~s.
        \\

        Primary material-age bound $A_{\max}$
        &
        \CoreMaterialAge
        &
        Measured from creation, not activation.
        Copy the final per-model setting; do not infer it from
        $C_{\max}$.
        \\

        Primary reuse-gap bound $G_{\max}$
        &
        \CoreReuseGap
        &
        Bound on the permitted gap since successful use.
        Retain any per-model differences explicitly.
        \\

        Primary payload profile
        &
        1000~B frames; 100~B/s
        &
        Complete alternating directional frames;
        10~s per frame.
        \\

        Service minimum
        &
        1000~B per direction per contact
        &
        Unchanged across primary load levels.
        \\

        Offered payload per direction
        &
        Low: 1000~B; medium: 3000~B
        &
        High demand follows Eq.~\eqref{eq:high-demand}.
        \\

        Management profile
        &
        Table~\ref{tab:costs}; 400~B/s
        &
        Synthetic operation durations, CPU, and message sizes.
        \\

        FixedPeriod interval
        &
        3600~s after successful SWITCH
        &
        One declared periodic baseline, not an optimized
        periodic-policy family.
        \\

        Matched horizons
        &
        Next: $H=1$; Multi: $H=2$
        &
        One additional $H=3$ sensitivity case is separately
        reported.
        \\

        Search limits
        &
        5000 nodes and 1~s
        &
        Limits apply per bounded search.
        \\

        Decision and process limits
        &
        2~s per decision; 600~s per process
        &
        No additional cumulative controller-time cutoff.
        \\

        Seed and suffix limits
        &
        350~ms and 10\,000 transitions each
        &
        Included in the common decision budget.
        \\

        Measured-delay parameters
        &
        $\delta_{\mathrm{fixed}}=0$, $\alpha=1$
        &
        Charges measured host decision time on the modeled
        critical path.
        \\

        Software and JVM
        &
        Java 21.0.10; Maven 3.9.9; 1~GiB heap
        &
        The same frozen final-v2 JAR is used throughout
        the campaign; JVMs run serially.
        \\

        Host CPU, RAM, and OS
        &
        \FinalHostHardware
        &
        Host timings are not onboard benchmarks.
        \\

        \bottomrule
    \end{tabularx}
\end{table*}

For a contact of duration $L_i$, high demand offers
\begin{equation}
    B_i^{\mathrm{high}}
    =
    1000
    \max\left(
        3,
        \left\lfloor
            \frac{L_i}{10\,\mathrm{s}}
        \right\rfloor
    \right)\mathrm{B}
    \label{eq:high-demand}
\end{equation}
in each direction.

The service threshold does not change with offered demand.
Thus a load intervention changes the amount available for
delivery without redefining a served contact.

The 54-case offline core compares FixedPeriod, just-in-time
PerContact, Next, Multi, and no-DFS seed at both horizons
for three missions and three load levels.
Primary geometric runs begin with empty key stores.

The 15 measured cases repeat the three ISS loads for the
first five configurations, omitting seed $H=2$.

Sensitivity is one factor at a time on medium ISS:
INSTALL phase durations are multiplied by 0.5 or 2 while
byte counts remain unchanged;
the cryptoperiod changes to 2700 or 10\,800~s;
and one Multi $H=3$ case uses the base model.
Duration scaling changes modeled processing and radio timing;
it is not hardware calibration.

X5, Y, and the previous three-contact high-demand witness
are executed under both tie rules offline and under
CANONICAL in measured mode.

Disclosure cases preserve the registered undetected,
two-second detection, gap detection, and late high-demand
incidents.
Channel cases preserve the original temporary X interruption
and add a preregistered two-second loss beginning one second
after the first ISS contact opens.
The actual interrupted operation is reported without moving
the outage to fit an observed policy decision.

Each registered cell receives one attempt without a budget
increase or automatic experiment retry.
Independent cells continue after recorded technical failures.
All incomplete and failed attempts remain visible.

The H3 factor is compared with its registered H2 baseline;
it is not counted as another H1/H2 pair.

\subsection{Input Structure and Reuse Opportunity}
\label{sec:input-structure}

Input-only analysis finds minimum selected-window lengths
of 58.98, 122.69, and 74.82~s for ISS, NOAA-19, and HST.
There are respectively one, five, and three adjacent
A--B--A triples.

All windows exceed the 26~s prepared-entry
INSTALL--SWITCH--OPEN--minimum-payload chain.
None lies in the 22--26~s band distinguishing that chain
from the already installed entry.

These are conditional duration checks, not achievable-service
guarantees.
They omit other validity and resource constraints and measured
decision delay.
Increasing offered traffic alone therefore does not reproduce
the narrow short-return feasibility condition in the primary
geometric core.
The separate X5 witness retains that mechanism test.

All 29 ISS, 22 NOAA-19, and 35 HST same-station recurrence
gaps exceed the base 5400~s cryptoperiod.
For a key activated at $a$ no later than the previous contact
end $e$, and return start $s$,
\begin{equation}
    d_k
    \leq a+5400
    \leq e+5400
    <s.
    \label{eq:no-active-reuse}
\end{equation}

Consequently, already activated material cannot survive to
the next same-station visit under the primary setting.
This does not prohibit an unactivated installed candidate,
whose activation clock has not started.
It also does not cover the 10\,800~s sensitivity.

The same premises give a directional-key-count bound.
Every served contact requires actual delivery in both
directions.
Key bindings prevent sharing across stations or directions,
while Eq.~\eqref{eq:no-active-reuse} prevents reuse at another
same-station contact.
With empty initial key stores,
\begin{equation}
    N_{\mathrm{new}}
    \geq
    2N_{\mathrm{served}}.
    \label{eq:key-bound}
\end{equation}

Earlier staged preparation still counts its generated keys.
Abandoned preparation can only add to the count.
This is a key-count bound only:
it proves neither CPU/byte lower bounds, useful-time optimality,
nor global policy optimality.
It does not extend to the longer-cryptoperiod sensitivity.

\subsection{Service, Security, and Computation Metrics}

Alongside the five objective components, we report full
execution status, search status, elapsed time, explored nodes,
and practical incumbent/fallback use.

Computational effort is not the modeled security-operation CPU.
A process failure, search limit, and failure to serve a contact
are distinct outcomes.

Process elapsed time includes startup and final export.
Controller elapsed time spans the controller call, including
checkpoint callbacks but excluding final serialization.
Aggregate thread CPU and garbage-collection samples are taken
later by the harness and therefore cover a different interval.
They are not interpreted as matched-interval CPU utilization.
Individual search-ledger measurements are reported separately.

Disclosure experiments reveal a selected active traffic-key
pair to a recording attacker.
Replacement material is assumed independent, and management
credentials remain uncompromised.
Detection notice is an exogenous event, not the output of
an evaluated detector.

Recovery is separated into:
spacecraft notice,
independent-pair activation,
independent-SA opening, and
first independent delivered frame.

A milestone not reached in a completed detected case is
right-censored, not assigned zero.
Undetected incidents and no-attack cases remain explicitly
inapplicable.
An incomplete execution has no final quality result.

Confidential service is retrospective:
each direction's delivered minimum must be met by complete
frames whose keys remain undisclosed at experiment end.
This excludes earlier traffic under a subsequently disclosed
key and never feeds future knowledge into planning.
It does not imply that every byte in the contact is confidential.

Emitted exposure includes aborted frame prefixes.
Exposure and confidential delivery are therefore not
complementary parts of one common delivered-byte total.

\subsection{Evidence and Reproduction}

The accompanying reproduction package retains exact scenario
models, source and executable identities, archived public
inputs, per-case invocations, journals, checkpoints, outcome
tables, and figure provenance.

The full case mapping is recorded in
\path{docs/final-research/runs.csv}.
The replay helper
\path{scripts/replay_final_case.py}
checks the frozen input and executable identities and applies
the original limits.
Replays must use a new output directory rather than overwrite
registered outcomes.

The separately compiled supplement contains complete outcome,
search, decision-delay, disclosure, recovery, and interruption
tables.
A repeated deterministic replay is a reproducibility check,
not an independent policy-effectiveness sample.

% AUTHOR CHECK:
% Add the permanent public repository or archive DOI here when
% it actually exists. Do not invent a deposit URL.

\FloatBarrier

% ============================================================
% VII. RESULTS
% ============================================================

\section{Results}
\label{sec:results}

Earlier SE2, R1, and PC2 artifacts remain a separate historical
record.
Their attempts, focused tests, and selected witnesses are not
counted as new observations in the final campaign.

Of the nine primary mission--load comparisons, four have
completed executions for both Next and Multi.
All four are equal in every component of
Eq.~\eqref{eq:objective}.
The remaining five pairs are incomplete because at least one
execution reaches the process watchdog.

The completed primary comparisons therefore establish no
practical quality advantage from the additional horizon at
the fixed budget.
This conclusion neither disproves the conditional exact
short-return mechanism nor assigns a ranking to unfinished
pairs.

Table~\ref{tab:core} reports the primary offline outcomes.
Among all 108 registered attempts, 93 reached
\op{FINISHED} and 15 ended with
\op{PROCESS\_TIMEOUT}.
The event-log audit checks all 93 exported executions and
37\,164 same-root full-comparator selections.

\begin{table*}[!t]
    \centering
    \small
    \setlength{\tabcolsep}{4pt}
    \renewcommand{\arraystretch}{1.15}

    \caption{
        Final offline core.
        Each numeric cell reports served contacts /
        useful secure transmission time in seconds.
        I denotes unavailable final quality after a process
        timeout, not zero service.
    }
    \label{tab:core}

    \begin{tabular*}{\textwidth}{
        @{\extracolsep{\fill}}
        llrrrrrr
        @{}
    }
        \toprule

        Mission
        & Load
        & FixedPeriod
        & PerContact
        & Next
        & Multi
        & Seed $H=1$
        & Seed $H=2$ \\

        \midrule

        ISS46
        & Low
        & $46/920$
        & $46/920$
        & $46/920$
        & $46/920$
        & $46/920$
        & $46/920$ \\

        ISS46
        & Medium
        & $46/2740$
        & $46/2740$
        & $46/2740$
        & \missing
        & $46/2740$
        & $46/2740$ \\

        ISS46
        & High
        & $46/14\,390$
        & $46/14\,390$
        & \missing
        & \missing
        & $46/14\,390$
        & $46/14\,390$ \\

        \addlinespace[3pt]

        NOAA38
        & Low
        & $38/760$
        & $38/760$
        & $38/760$
        & $38/760$
        & $38/760$
        & $38/760$ \\

        NOAA38
        & Medium
        & $38/2280$
        & $38/2280$
        & $38/2280$
        & $38/2280$
        & $38/2280$
        & $38/2280$ \\

        NOAA38
        & High
        & $38/19\,210$
        & $38/19\,210$
        & \missing
        & \missing
        & $38/19\,210$
        & $38/19\,210$ \\

        \addlinespace[3pt]

        HST48
        & Low
        & $48/960$
        & $48/960$
        & $48/960$
        & $48/960$
        & $48/960$
        & $48/960$ \\

        HST48
        & Medium
        & $48/2880$
        & $48/2880$
        & $48/2880$
        & \missing
        & $48/2880$
        & $48/2880$ \\

        HST48
        & High
        & $48/15\,560$
        & $48/15\,560$
        & \missing
        & \missing
        & $48/15\,560$
        & $48/15\,560$ \\

        \bottomrule
    \end{tabular*}

    \vspace{4pt}

    \begin{minipage}{0.98\textwidth}
        \footnotesize
        \emph{Interpretation:}
        Next uses $H=1$ and Multi uses $H=2$.
        The slash separates two measurements; it is not a ratio.
        Eight executions are unfinished in this table, producing
        five incomplete Next/Multi pairs.
        Equality of the four completed pairs applies to all five
        objective components, not only the two displayed here.
        Full status and cost tables are retained in the supplement.
    \end{minipage}
\end{table*}

\subsection{Prepared Mechanism versus Bounded Autonomous Execution}

The prepared exact X witness remains a validated conditional
mechanism:
Next serves two contacts with 38~s of useful transmission,
whereas Multi serves three with 37~s by installing A material
before the short return.

After the prescribed common PREPARE, the reference uses the
automatic exact planner and receding controller, with every
search asserted \op{COMPLETE}.
The preliminary preparation contributes 0.20~s of operation CPU
and two directional keys exactly once.

The full-suite reference permits 2\,000\,000 nodes per decision
without a search elapsed-time deadline.
It is not a result of the final 5000-node autonomous practical
configuration.
Completed visible-horizon searches do not establish global
whole-mission optimality.

In the final autonomous X5 regressions, both practical horizons,
under both tie rules, serve two contacts with 38~s useful time,
0.12~s modeled operation CPU, 160 management bytes,
and no newly generated keys.

At the decisive $t=0$ Multi root, all three contacts are visible,
but search stops at 5000 nodes.
The available seed and retained candidates serve two contacts
and prepare only near the later return.
The search candidate serves the same two contacts more cheaply.

No available three-contact candidate is rejected by the portfolio
comparator.
This is an observed bounded-search limitation, not proof that
a three-contact continuation is infeasible.

Both horizons serve all three Y contacts with 42~s offline.
On the retained high-demand witness, all offline tie-rule
variants and both measured variants serve three contacts with
410~s, 2.13~s operation CPU, 3240 management bytes,
and six new directional keys.

In measured X5, Next and Multi still serve two contacts,
but useful time becomes 33.75 and 27.5~s, respectively.
Their accumulated measured decision delays are approximately
4.617 and 11.392~s, with no deadline misses.

Equal nominal costs therefore do not imply equal executed
payload when computation occupies the critical path.
These are individual registered executions, not statistical
latency estimates.

\subsection{Implementation Verification}

A fresh isolated Maven clean build passes all 478 tests in
64 classes, with no failures, errors, or skips.
The first attempt's five local HTTP errors and the successful
rebuild have identical 181-file source manifests.
Local loopback permission, rather than a source change or
test selection, changed between the builds.
Both records are preserved.

The same shaded JAR supplies every final offline and measured
invocation.
The supplement records its SHA-256 and source identity.
A base Git commit alone is insufficient because the working
source contained explicit changes.

The full suite covers checkpoint restoration, lifecycle and
counter constraints, the prepared exact witness, and the
practical-portfolio tests.
Both small PC2 oracle searches are \op{COMPLETE} at
1364 nodes.

These checks cover their specified fixtures.
They do not establish universal bounded-search optimality,
cryptographic security, or statistical generalization.
Historical focused test counts overlap and are not added
to the 478-test total.

\subsection{Primary Geometric Workloads}

All 36 FixedPeriod, PerContact, and no-DFS seed executions
finish.
Six of nine Next executions and four of nine Multi executions
finish.
Every finished primary execution serves all selected contacts.

The four completed horizon pairs are low-load ISS,
low- and medium-load NOAA-19, and low-load HST.
They tie on all five objective components.
The five remaining pairs have unavailable final comparisons.

All 27 PerContact/seed executions attain the key-count bound
in Eq.~\eqref{eq:key-bound}.
This does not establish optimality of their other components.

For ISS, NOAA-19, and HST, respectively, finished
PerContact/Next/Multi/seed configurations use
32.66, 26.98, and 34.08~s of modeled operation CPU;
49\,680, 41\,040, and 51\,840 management bytes; and
92, 76, and 96 directional keys.

Raising offered traffic does not change these lifecycle totals
in the completed primary runs.

FixedPeriod produces 15, 14, and 11 additional never-activated
pairs for ISS, NOAA-19, and HST.
These add 3.00, 2.80, and 2.20~s of PREPARE CPU and
30, 28, and 22 directional keys.

FixedPeriod delivers the same service and useful time as the
other finished policies.
This is a cost disadvantage of the declared periodic rule,
not a loss of contact service or a claim about every possible
period.

Figures~\ref{fig:service} and \ref{fig:cost} show the primary
service and lifecycle-cost outcomes cell by cell, marking
timed-out executions explicitly instead of omitting them.

\widefigure
    {service_matrix.pdf}
    {0.48\textheight}
    {Primary offline service.
     Each cell gives served/selected contacts and useful secure
     transmission time for one configuration, mission, and load.
     Hatched cells reached the 600~s process watchdog; their final
     quality is unavailable, not zero.
     Every completed execution serves all selected contacts;
     equality of completed cells does not rank unfinished pairs.}
    {fig:service}

\widefigure
    {cost_matrix.pdf}
    {0.66\textheight}
    {Primary lifecycle costs, excluding search and payload-encryption
     CPU.
     Shading gives the overhead relative to the cheapest completed
     configuration in the same mission--load column.
     The declared FixedPeriod rule generates 23--37\% more directional
     keys and 6--10\% more modeled operation CPU from unused material;
     management bytes are identical among completed configurations.
     Hatched cells are process timeouts without final values.}
    {fig:cost}

The additional search does not improve any completed primary
horizon pair.
Low-load ISS Next/Multi use 213.28/521.23~s of controller
elapsed time, versus 16.29/23.17~s for seed H1/H2.

For low-load NOAA-19, the corresponding times are
139.61/389.14~s versus 11.53/15.84~s.
For medium-load NOAA-19, they are 337.17/569.84~s versus
17.86/24.92~s.
For low-load HST, they are 97.57/453.39~s versus
15.46/21.58~s.

These are individual host observations, not onboard benchmarks.
In low-load ISS, Next/Multi contain 209/593 incomplete
searches out of 570/614 attempted searches, yet both physical
executions finish because the practical portfolio admits
feasible uncertified plans.

Eight primary executions reach the 600~s process watchdog:
Multi on medium-load ISS and HST, and both horizons on
every high-load workload.
Figure~\ref{fig:runtime} reports process elapsed time, which
includes startup and export and differs from controller time.

Last saved prefixes remain available.
Neither their contact counts nor absent final search ledgers
are interpreted as completed quality.

All 93 finished campaign executions have zero decision
deadline misses.
This does not prevent cumulative computation from exceeding
the separate process budget.

Transient checkpoint replacement refusals are recorded and
retried.
The outside cause of Windows access denial is unknown.
Final exported snapshots agree with their journals and
checkpoints.

\widefigure
    {runtime_completion.pdf}
    {0.34\textheight}
    {Elapsed process time for each primary attempt on a logarithmic
     scale shared across missions.
     Lines connect the loads of one configuration; X markers on the
     registered 600~s watchdog line denote unfinished executions,
     joined to the last finished load by a dotted segment.
     These single host observations are not statistical or onboard
     timing estimates.}
    {fig:runtime}

\subsection{Measured Computation and Sensitivity}

Twelve of fifteen measured ISS executions finish.
Medium-load Multi and both high-load search policies reach
the process limit.

Low-load service and all five objective components equal
their offline counterparts despite 238.54/537.02~s of accumulated
measured decision delay for Next/Multi.

The medium-load Next execution retains 46 served contacts,
but useful time decreases from 2740 to 2720~s.
High-load seed H1 similarly decreases from 14\,390 to
14\,290~s while retaining 46 contacts.
FixedPeriod and PerContact retain their offline useful times
at all loads.

Figure~\ref{fig:measured} shows these changes relative to
offline execution and distinguishes them from incomplete cells.
Decision delay can consume payload opportunity without
violating an individual decision deadline.

Of nine sensitivity attempts, five finish.
Halving INSTALL phase durations gives Next 46 contacts,
2740~s useful time, and 25.76~s operation CPU.
Doubling them gives 46 contacts, 2730~s, and 46.46~s CPU.

Management bytes and key counts remain 49\,680 and 92.
The CPU change follows the explicitly scaled CPU phases
as well as elapsed/radio phases; byte counts were not scaled.

Both paired Multi executions time out.
These completed Next outcomes therefore do not establish
horizon advantages.

The 2700~s cryptoperiod leaves the completed Next five-component
score unchanged, while Multi again times out.
The H3 execution also reaches its watchdog and cannot be
ranked against the unfinished base H2 execution.

At a 10\,800~s cryptoperiod, both horizons finish and tie at
\[
    (46,\;2740,\;22.91,\;34\,680,\;62).
\]

Their journals use 31 key pairs:
15 pairs in two contacts and 16 pairs in one contact.
Relative to completed base Next, each saves 15
PREPARE/INSTALL/SWITCH cycles, 9.75~s operation CPU,
15\,000 management bytes, and 30 directional keys.
There remain 46 OPEN operations.

This is observed material reuse enabled by a longer
cryptoperiod, not a unique Multi advantage.
The two plans first differ in preparation timing, but reach
the same installation time and final metrics.

Neither the base key-count bound nor a security conclusion
is extended to this longer-lifetime sensitivity.

\widefigure
    {measured_iss.pdf}
    {0.45\textheight}
    {Effect of charging measured host decision delay on ISS service.
     Bars give the change in useful secure transmission time
     (measured minus offline); the right columns list
     offline~$\rightarrow$~measured served contacts and useful time.
     Only medium-load Next ($-20$~s) and high-load seed $H=1$
     ($-100$~s) change, and served contacts are unchanged wherever
     both executions finish.
     Hatched rows are process timeouts in both modes.
     The modeled delay excludes checkpoint I/O and is not a measured
     flight deployment.}
    {fig:measured}

\subsection{Disclosure, Recovery, and Channel Loss}

All ten security/channel executions and both
management-interruption executions finish.

Figure~\ref{fig:exposure} separates retrospectively decryptable
emissions before disclosure from emissions afterward.

In undetected X, both horizons expose 7600~B and deliver
7600 confidential bytes.
With the original two-second detection, both expose 3600~B
and deliver 9800 confidential bytes.
These values include recorded past traffic under the later
disclosed key; they are not exposure only after detection.

For detected X, both horizons deliver spacecraft notice at
0.4~s, activate independent keys at 7.4~s, open an independent
SA at 8.4~s, and deliver the first independent frame at
8.65~s after ground notice.

Gap detection delivers notice only after 18.4~s.
Independent activation, SA opening, and frame delivery remain
censored at experiment end.
Delivering notification therefore does not establish completed
recovery.

First readiness is not continuous readiness.
The audit retains subsequent expiry and closure intervals.
Figure~\ref{fig:recovery} preserves the distinction between
these milestones and missing events.

The late high-demand disclosure produces the only completed
final pair with a Multi objective gain.
Both policies serve three contacts and use 2.88~s operation
CPU, 4416 management bytes, and eight keys.
Multi delivers 410 rather than 400 useful seconds.

Its confidential delivery is 7000 rather than 6000~B,
corresponding to one additional 1000-byte uplink frame.
Next/Multi emitted exposure is 34\,042/34\,000~B.
The extra 42~B for Next belong to an aborted frame prefix.
Emitted exposure must not be subtracted from complete delivery
as if the two populations were identical.

The journal identifies the timing mechanism.
Multi completes independent preparation before disclosure
and begins INSTALL before ground notice.
It finishes that operation before notification.

Spacecraft notice is later for Multi:
3.9753 rather than 0.4~s after ground notice.
However, independent activation occurs at 4.9753 rather than
7.4~s, SA opening at 5.9753 rather than 8.4~s, and first
independent delivery at 15.9753 rather than 18.4~s.

This is a consequence of nominal planning interacting with
the incident and available traffic.
The planner did not know the future attack or optimize recovery.

Offline controller times are 18.71~s for Multi and 3.35~s
for Next.
The gain is not a measured-latency security result or proof
of general recovery superiority.

The original temporary X channel loss intersects no in-flight
frame or online operation in either final execution.
Both finish at
\[
    (3,\;32,\;0.89,\;1320,\;2).
\]
This case does not validate mid-operation interruption.

The separate preregistered ISS3 outage does interrupt INSTALL
under both horizons.
Each execution records 0.10~s CPU and 360 uplink bytes
of failed work, retries through normal admission, and serves
all three contacts.

The final score is
\[
    (3,\;60,\;2.23,\;3600,\;6),
\]
with no aborted frame.

This directly checks consumed partial work and INSTALL retry.
It does not exercise reconciliation after uncertain SWITCH
or OPEN, which remains unimplemented.

\widefigure
    {detected_exposure.pdf}
    {0.33\textheight}
    {Decryptable emitted traffic under the recording-attacker model:
     bytes recorded before disclosure, bytes emitted after
     disclosure, and their combined share of all emitted bytes,
     with color identifying policy.
     The channel-only case has zero exposure because no disclosure
     is injected.
     Exposure includes aborted prefixes and differs from delivered
     traffic.}
    {fig:exposure}

\widefigure
    {recovery.pdf}
    {0.42\textheight}
    {Four recovery milestones measured from ground-station notice
     for the detected disclosure cases, one lane per milestone.
     Dashed arrows mark milestones censored at experiment end,
     not zero delay.
     First activation or SA readiness is not a continuous-readiness
     guarantee.
     In the late-demand case, Multi receives spacecraft notice later
     (3.98 vs.\ 0.4~s) but delivers its first independent frame
     earlier (15.98 vs.\ 18.4~s).}
    {fig:recovery}

% Flush all results figures before Discussion and References.
\FloatBarrier

% ============================================================
% VIII. DISCUSSION
% ============================================================

\section{Discussion and Limitations}
\label{sec:discussion}

\subsection{What the Evidence Establishes}

The answer to the research question has three parts.

First, additional visibility has a concrete finite-model
service mechanism when earlier installation changes the
feasibility of a short return.
The prepared exact witness verifies that mechanism.

Second, the final bounded autonomous configuration does not
necessarily find it.
In X5, the portfolio never receives the known three-contact
candidate before the search limit.

Third, the geometric core supplies no completed quality
advantage from the extra horizon.
Every seed control finishes, while search frequently exhausts
the process budget.

These observations provide no practical advantage for additional
horizon under the tested primary configuration.
They do not prove that contact prediction is universally
useless or that unfinished pairs would tie.

\subsection{Why the Primary Workloads Offer Limited Opportunity}

The input structure explains part of the observed equality
without attributing motives to the optimizer.

Every primary contact window exceeds the specified
prepared-entry minimum.
No already activated key can survive a same-station return
under the base cryptoperiod.

Increasing offered traffic alone does not create the short-return
feasibility contrast in Eq.~\eqref{eq:opportunity}.
Extending the cryptoperiod creates actual reuse, but both
horizons realize the same savings.

A longer permitted lifetime is not itself a security
improvement.
It changes potential exposure as well as management work.

\subsection{Security Consequences and Conditional Benefits}

The late-disclosure witness demonstrates a separate conditional
effect:
nominal preparation and installation timing can improve
independent delivery during a later incident.

However, spacecraft notification is slower, and much of the
previous traffic remains retrospectively exposed.
Recovery delay, useful time, and disclosure exposure are
different outcomes.

Across the completed registered horizon pairs, there are
19 equal outcomes, one Next advantage in measured X5,
and one Multi advantage in late disclosure.
Ten pairs remain incomplete, and H3 is separately unpaired.

These heterogeneous cases are not independent samples from
which to infer a population win rate.
They support reporting mechanisms, conditions, and costs,
not a universal policy ranking.

\subsection{Computational and External-Validity Limits}

The exact reference is an offline finite-model benchmark.
The practical mode trades certification for bounded decision
effort and must be evaluated as a different controller.
A completed physical execution does not imply completed
searches at every decision.

Delay experiments expose selected critical-path assumptions.
They do not measure a full ground/onboard deployment.
Controlled losses test particular interventions, not a general
distribution of contact-prediction errors or detector behavior.

Operation and traffic costs remain synthetic.
The CPU objective includes lifecycle-management operations,
not per-byte payload encryption, energy, or a complete flight
processor load.

Public propagated visibility provides geometric realism,
not operational access, link budgets, or authorization for
actual ISS, NOAA-19, or HST communications.
The workloads are not independent samples of operational
missions.
Repeated deterministic runs are not independent
policy-effectiveness observations.

An incomplete execution cannot support a claim that Multi
is better, equal, or never worse.

\subsection{Security-Model Boundaries}

The threat model excludes persistent key-generator compromise,
active forgery, a physical RF jammer model, compromised
management credentials, and a real distributed endpoint
protocol.

The nonce-accounting correction addresses one concrete
simulation invariant.
It does not implement encryption or prove replay resistance.

Explicit reconciliation after uncertain online operations
remains absent.
These are limitations of the current model and evidence,
not completed features.

% ============================================================
% IX. CONCLUSION
% ============================================================

\section{Conclusion}
\label{sec:conclusion}

Joint traffic-key and SA lifecycle timing admits a conditional
multi-contact scheduling mechanism, but observing more
contacts is insufficient to realize it efficiently.

In one frozen 108-attempt campaign, the four completed
geometric-core horizon pairs tie on all five objective
components, and five pairs remain incomplete.
No-DFS controls finish all primary workloads.

The prepared exact witness retains an additional served
contact, while the bounded autonomous witness does not.
A late-disclosure case yields an extra confidential frame
under Multi.
Conversely, charging host decision time reduces Multi's
payload in measured X5.

INSTALL interruption is handled with consumed partial costs
and an admissible retry in the tested cases.
These results establish conditional opportunities and explicit
computational limits, not broad Multi superiority.

The contribution is the controlled lifecycle formulation,
shared executor, local available-candidate guarantee,
and reproducible evidence separating forecast depth,
feasibility, computation, and recovery.

Operational validation, concrete distributed reconciliation,
calibrated mission costs, and cryptographic-security proof
remain outside this study.

\FloatBarrier

% ============================================================
% REFERENCES
%
% Ordered by first citation:
% SDLS, EP, SKM, IOTA, BPv7, BPSec, SABR, Detke,
% receding-horizon control, CelesTrak, SatNOGS, Orekit.
% ============================================================

\begin{thebibliography}{12}

\bibitem{sdls}
Consultative Committee for Space Data Systems,
\emph{Space Data Link Security Protocol},
CCSDS 355.0-B-2,
July 2022.
[Online]. Available:
\url{https://ccsds.org/Pubs/355x0b2.pdf}.

\bibitem{ep}
Consultative Committee for Space Data Systems,
\emph{Space Data Link Security Protocol---Extended Procedures},
CCSDS 355.1-B-1,
Feb. 2020.
[Online]. Available:
\url{https://ccsds.org/Pubs/355x1b1.pdf}.

\bibitem{skm}
Consultative Committee for Space Data Systems,
\emph{Symmetric Key Management},
CCSDS 354.0-M-1,
Dec. 2023.
[Online]. Available:
\url{https://ccsds.org/Pubs/354x0m1.pdf}.

\bibitem{iota}
D. Bernal, O. Ledesma, P. Lamo, and J. Bermejo,
``IOTA-Based Authentication for IoT Devices in Satellite Networks,''
\emph{Computers, Materials \& Continua},
vol. 86, no. 1, pp. 1--39, 2026,
doi: 10.32604/cmc.2025.069746.

\bibitem{bp}
S. Burleigh, K. Fall, and E. Birrane III,
``Bundle Protocol Version 7,''
RFC 9171,
Jan. 2022.
[Online]. Available:
\url{https://www.rfc-editor.org/rfc/rfc9171.html}.

\bibitem{bpsec}
E. Birrane III and K. McKeever,
``Bundle Protocol Security (BPSec),''
RFC 9172,
Jan. 2022.
[Online]. Available:
\url{https://www.rfc-editor.org/rfc/rfc9172.html}.

\bibitem{sabr}
Consultative Committee for Space Data Systems,
\emph{Schedule-Aware Bundle Routing},
CCSDS 734.3-B-1,
July 2019.
[Online]. Available:
\url{https://ccsds.org/Pubs/734x3b1.pdf}.

\bibitem{detke}
R. Detke, P. G. Madoery, R. Cherini, J. A. Fraire,
and J. M. Finochietto,
``Security-Aware Route Selection in DTNs:
Towards Optimal Key Provisioning,''
IEEE WiSEE,
Space-Terrestrial Internetworking Workshop (STINT), 2026.
Conference-hosted full paper.
Accessed: Oct. 25, 2026.
[Online]. Available:
\url{https://stintworkshops.org/papers/stint2026-19-detke.pdf}.

\bibitem{rhc}
J. Mattingley, Y. Wang, and S. Boyd,
``Code Generation for Receding Horizon Control,''
in \emph{Proc. IEEE International Symposium on
Computer-Aided Control System Design},
Yokohama, Japan,
2010, pp. 985--992.
[Online]. Available:
\url{https://stanford.edu/~boyd/papers/pdf/code_gen_rhc_msc.pdf}.

\bibitem{celestrak}
T. S. Kelso,
``A New Way to Obtain GP Data (aka TLEs),''
CelesTrak.
Accessed: Oct. 25, 2026.
[Online]. Available:
\url{https://celestrak.org/NORAD/documentation/gp-data-formats.php}.

\bibitem{satnogs}
Libre Space Foundation,
``SatNOGS Network API Documentation.''
Accessed: Oct. 25, 2026.
[Online]. Available:
\url{https://docs.satnogs.org/projects/satnogs-network/en/stable/api.html}.

\bibitem{orekit}
Orekit,
``Orekit: An Open-Source Library for Spaceflight Dynamics.''
Accessed: Oct. 25, 2026.
[Online]. Available:
\url{https://www.orekit.org/}.

\end{thebibliography}

\end{document}
